\writeSectionFrame{\presentationTitle}{Fazit}
\begin{frame}{Immutable objects}
	\begin{itemize}
	  \item Immutable Objects verschieben Komplexität
	  \item Weg von verstecktem Zustand
	  \item Hin zu expliziten Ergebnissen
	  \item Unveränderlichkeit erzwingt Design, nicht nur Syntax!
	\end{itemize}	
\end{frame}

\begin{frame}{Vor- und Nachteile}
	\begin{exampleblock}{Vorteile}
		\begin{itemize}
		  \item Keine versteckten Seiteneffekte
		  \item Objekte sind immer in einem gültigen Zustand
		  \item Sehr gute Testbarkeit
		  \item Nebenläufigkeit ohne Synchronisation
		\end{itemize}	
	\end{exampleblock}
\end{frame}

\begin{frame}{Vor- und Nachteile}
	\begin{exampleblock}{Nachteile}
		\begin{itemize}
		  \item Mehr Objekte und mehr Rückgabewerte
		  \item Mehr Speicherbedarf und geringere Performanz
		  \item Änderungen wirken manchmal umständlicher
		  \item Nicht jeder Zustand ist sinnvoll immutable
		  \item Erfordert Umdenken im Design
		\end{itemize}
	\end{exampleblock}
\end{frame}

\begin{frame}{Wann sind Immutable Objects sinnvoll?}
	\begin{itemize}
	  \item Value Objects und DTOs
	  \item Domain-Modelle mit klaren Zustandsübergängen
	  \item Konfigurationsobjekte
	  \item Parallel oder nebenläufig genutzte Objekte
	\end{itemize}
\end{frame}

\begin{frame}{Wann eher nicht?}
	\begin{itemize}
	  \item Sehr große, häufig veränderte Datenstrukturen
	  \item Performance-kritische Low-Level-Strukturen
	  \item Kurzlebige technische Hilfsobjekte
	\end{itemize}
\end{frame}

\begin{frame}{Fazit}
	\begin{block}{}
		Immutable Objects machen Zustandsänderungen explizit. \\
		Das erhöht zunächst den Aufwand, \\
		reduziert aber langfristig die Komplexität.
	\end{block}
\end{frame}

